\documentclass[conference]{IEEEtran}

\usepackage[T1]{fontenc}
\usepackage[utf8]{inputenc}
\usepackage{amsmath}
\usepackage{booktabs}
\usepackage{cite}
\usepackage{graphicx}
\usepackage{hyperref}
\usepackage{pgfplots}
\usepackage{tikz}

\pgfplotsset{compat=1.18}

\newcommand{\MaybeInput}[2]{%
  \IfFileExists{#1}{\input{#1}}{#2}%
}

\title{Engineering Trade-offs in Huffman Compression:\\
An Experimental Comparison of Swift, Rust, and Streaming Implementations}

\author{
\IEEEauthorblockN{Hz Project}
\IEEEauthorblockA{Reference Swift implementation and benchmark suite}
}

\begin{document}

\maketitle

\begin{abstract}
The hz repository contains multiple implementations of the same Huffman archive format: a Swift reference engine, a Rust native backend called through a C ABI, and Rust bounded-memory file streaming. This paper studies implementation-level trade-offs among those execution modes without treating the experiment as a general language-performance contest. The planned measurements compare throughput, archive ratio, decompression time, peak resident memory, workload distribution, input-size scaling, and Swift--Rust integration overhead where it can be isolated. Results are generated by repository scripts and inserted into the paper as reproducible tables and figures; unsupported benchmark values are intentionally absent from this draft.
\end{abstract}

\section{Introduction}

Huffman coding is a well-understood entropy coding method that constructs a prefix code from observed symbol frequencies \cite{huffman1952method,sayood2017introduction}. The algorithm is not new, but its implementation context matters. In hz, the same archive format is implemented in Swift and Rust, and the Rust backend exposes both in-memory and streaming file paths. This makes the repository a useful setting for studying engineering effects that are often mixed together: algorithmic behavior, programming language runtime behavior, memory ownership, file I/O strategy, and foreign-function interface overhead.

This paper follows the recursive Huffman study in the same repository \cite{hzRecursiveHuffman}. The earlier paper asks what happens when Huffman archives are recursively compressed. This paper asks a different question: how do the current implementation strategies compare when they encode the same single-layer archive format under controlled workloads?

The comparison is deliberately narrow. It does not claim that Swift or Rust is universally faster, or that streaming is universally better than in-memory processing. Instead, it measures the actual hz execution modes under reproducible conditions and records enough metadata to separate algorithmic effects from implementation effects.

\section{Background}

The `.hz' archive format stores the metadata needed to reconstruct a deterministic Huffman tree: original byte count, encoded bit count, and a sorted byte frequency table. The encoded payload is packed most-significant-bit first. The encoded bit count and original byte count prevent padding bits from becoming decoded output.

The current version 2 archive stores frequencies rather than serialized canonical code lengths. The implementations may derive canonical code tables internally, but compatible decoders reconstruct the tree-path codes from the frequency table. This matters for benchmarking because all engines are expected to produce compatible archives for the same input, while differences in execution time and memory come from implementation strategy rather than a different compressed representation.

\section{Implementations Under Study}

The experiment compares three execution modes:

\begin{itemize}
    \item \textbf{Swift in-memory}: \texttt{SwiftHuffmanEngine} reads input into Swift \texttt{Data}, builds an archive in memory, and decompresses from an in-memory archive.
    \item \textbf{Rust in-memory}: \texttt{RustHuffmanEngine} crosses the Swift--Rust C ABI with complete buffers and calls the Rust single-layer codec.
    \item \textbf{Rust streaming}: the Rust path-based file API reads and writes through bounded buffers and temporary destination files.
\end{itemize}

\MaybeInput{tables/implementation-summary.tex}{%
\begin{table}[t]
\caption{Implementation summary placeholder. Run \texttt{./artifacts/reproduce.sh --smoke} to generate the table.}
\label{tab:implementation-summary}
\centering
\begin{tabular}{llll}
\toprule
Mode & Language & Memory model & Archive layers \\
\midrule
Swift & Swift & In-memory & Single-layer in this study \\
Rust & Rust via C ABI & In-memory buffers & Single-layer \\
Rust streaming & Rust via C ABI & Bounded file streaming & Single-layer \\
\bottomrule
\end{tabular}
\end{table}
}

Recursive and adaptive compression are not part of the primary implementation comparison. They are treated as a separate experiment because comparing recursive Swift output with single-layer Rust streaming output would confound compression mode with implementation mode.

\section{Experimental Methodology}

The primary experiment uses single-layer compression for every engine:

\begin{verbatim}
--max-depth 0
\end{verbatim}

Each engine processes the same generated input files. The benchmark runner verifies every produced archive by decompression and byte comparison. The reproduction script normalizes output into a stable CSV schema and records the compression mode, requested maximum recursion depth, accepted layer count, verification status, and process exit status.

Peak resident memory is captured on macOS with \texttt{/usr/bin/time -l} when available. The value is recorded as peak RSS in bytes. This measurement is process-level and includes the benchmark invocation environment; the limitation is reported with the data rather than hidden.

The implementation currently records wall-clock compression and decompression time from the benchmark runner. Throughput is computed as original input bytes divided by the corresponding elapsed time. Archive ratio is computed as archive bytes divided by original bytes.

\section{Workloads}

The workload generator creates deterministic files with controlled byte distributions. The full campaign is designed to include:

\begin{itemize}
    \item single-symbol data;
    \item balanced two-symbol data;
    \item skewed binary categorical data;
    \item uniform distributions over small alphabets;
    \item uniform distributions over all 256 byte values;
    \item Zipf-like categorical distributions;
    \item Markov-correlated byte sequences;
    \item ordinary prose;
    \item source-code-like text;
    \item high-entropy pseudorandom bytes;
    \item compressed-like deterministic bytes;
    \item existing recursive-compression-style repetitive workloads.
\end{itemize}

Each generated workload records its distribution family, seed, input size, alphabet size when applicable, and empirical byte entropy. Entropy is computed from the generated bytes rather than asserted from the intended distribution.

\MaybeInput{tables/workload-definitions.tex}{%
\begin{table}[t]
\caption{Workload definitions placeholder. Run the reproduction script to generate this table from the workload manifest.}
\label{tab:workloads}
\centering
\begin{tabular}{lrr}
\toprule
Workload & Distribution & Bytes \\
\midrule
TODO(results) & generated by scripts & -- \\
\bottomrule
\end{tabular}
\end{table}
}

\section{Metrics}

The normalized benchmark data records:

\begin{itemize}
    \item workload name, distribution family, seed, input size, and empirical byte entropy;
    \item engine and compression mode;
    \item requested maximum recursion depth and accepted recursive layers;
    \item original size, archive size, and compression ratio;
    \item compression and decompression wall time;
    \item compression and decompression throughput;
    \item verification status and process exit status;
    \item peak RSS where supported by the platform tooling.
\end{itemize}

These metrics support the following research questions:

\begin{description}
    \item[RQ1] How do Swift in-memory, Rust in-memory, and Rust streaming implementations compare in compression and decompression throughput for single-layer Huffman archives?
    \item[RQ2] How does peak RSS scale with input size for each execution model when measured with the same process-level mechanism?
    \item[RQ3] How do controlled byte distributions affect archive ratio and runtime?
    \item[RQ4] What overhead is observable when the Rust in-memory backend is called through the Swift--Rust integration boundary?
    \item[RQ5] When does bounded-memory streaming become preferable to in-memory processing for the hz file workflow?
\end{description}

\section{Results}

% TODO(results): generated by papers/implementation-benchmarks/artifacts/reproduce.sh
% The paper intentionally contains no invented benchmark numbers.

\MaybeInput{tables/environment.tex}{%
\begin{table}[t]
\caption{Benchmark environment placeholder. Run the reproduction script to capture the local environment.}
\label{tab:environment}
\centering
\begin{tabular}{ll}
\toprule
Field & Value \\
\midrule
TODO(results) & generated by reproduce.sh \\
\bottomrule
\end{tabular}
\end{table}
}

\MaybeInput{tables/primary-results.tex}{%
\begin{table*}[t]
\caption{Primary results placeholder. Run the reproduction script to generate verified benchmark rows.}
\label{tab:primary-results}
\centering
\begin{tabular}{llll}
\toprule
Workload & Engine & Ratio & Verified \\
\midrule
TODO(results) & generated by reproduce.sh & -- & -- \\
\bottomrule
\end{tabular}
\end{table*}
}

\begin{figure}[t]
\centering
\MaybeInput{figures/compression-throughput.tex}{\fbox{\parbox{0.88\columnwidth}{\centering TODO(results): compression throughput figure generated by reproduce.sh}}}
\caption{Compression throughput by engine and workload.}
\label{fig:compression-throughput}
\end{figure}

\begin{figure}[t]
\centering
\MaybeInput{figures/decompression-throughput.tex}{\fbox{\parbox{0.88\columnwidth}{\centering TODO(results): decompression throughput figure generated by reproduce.sh}}}
\caption{Decompression throughput by engine and workload.}
\label{fig:decompression-throughput}
\end{figure}

\begin{figure}[t]
\centering
\MaybeInput{figures/ratio-vs-entropy.tex}{\fbox{\parbox{0.88\columnwidth}{\centering TODO(results): archive ratio versus empirical entropy generated by reproduce.sh}}}
\caption{Archive ratio versus empirical byte entropy.}
\label{fig:ratio-entropy}
\end{figure}

\begin{figure}[t]
\centering
\MaybeInput{figures/peak-rss-vs-input-size.tex}{\fbox{\parbox{0.88\columnwidth}{\centering TODO(results): peak RSS versus input size generated by reproduce.sh}}}
\caption{Peak RSS versus input size.}
\label{fig:rss-input-size}
\end{figure}

\begin{figure}[t]
\centering
\MaybeInput{figures/rust-stream-overhead.tex}{\fbox{\parbox{0.88\columnwidth}{\centering TODO(results): Rust streaming overhead generated by reproduce.sh}}}
\caption{Rust streaming compression time delta relative to Rust in-memory mode.}
\label{fig:stream-overhead}
\end{figure}

\section{Discussion}

The planned analysis distinguishes four sources of variation. First, byte distribution affects the Huffman tree, archive ratio, and encoded bit count. Second, implementation language affects data structures, allocation behavior, and bit I/O cost. Third, memory model affects peak RSS: in-memory engines operate on complete buffers, while streaming reads and writes bounded chunks. Fourth, FFI integration affects Rust in-memory calls because input and output cross the Swift--Rust boundary as complete buffers.

The Rust streaming path is not expected to dominate every timing metric. It performs file I/O and uses a two-pass compressor over seekable input. Its intended advantage is bounded memory for file workflows, not necessarily lower latency on small inputs. The experiment therefore reports both time and memory rather than treating throughput as the only outcome.

\section{Threats to Validity}

The benchmark workloads are generated rather than drawn from a large public corpus. This improves reproducibility and distribution control but limits external validity. The benchmark runner measures the current hz implementation rather than optimized production compressors. Peak RSS is collected at the process level and can include invocation overhead. Smoke runs are intended for validation only and should not be interpreted as final experimental evidence.

The Rust in-memory engine is called through Swift in the benchmark runner. That is appropriate for measuring the integrated application path, but it does not isolate pure Rust library performance. A future native-only harness could separate C ABI overhead from Rust codec cost more precisely.

\section{Reproducibility}

The smoke campaign can be run with:

\begin{verbatim}
papers/implementation-benchmarks/artifacts/reproduce.sh --smoke
\end{verbatim}

The full campaign can be run with:

\begin{verbatim}
papers/implementation-benchmarks/artifacts/reproduce.sh
\end{verbatim}

The script generates workloads, runs each engine, verifies decompression, records normalized CSV data, captures environment metadata, and generates paper-ready TeX fragments under \texttt{tables/} and \texttt{figures/}. Generated raw outputs, logs, workloads, local environment captures, figures, and result CSVs are ignored by default.

\section{Related Work}

Huffman's original paper introduced the minimum-redundancy prefix-code construction \cite{huffman1952method}. General compression texts discuss Huffman coding as part of broader source coding and compression systems \cite{sayood2017introduction}. The companion hz recursive-compression manuscript studies recursive archive behavior and provides background on the repository's recursive mode \cite{hzRecursiveHuffman}.

\section{Conclusion}

This draft establishes the structure for a reproducible implementation benchmark study of hz. The central comparison is not Swift versus Rust in the abstract, but rather the concrete trade-offs among Swift in-memory, Rust in-memory through a C ABI, and Rust bounded-memory streaming while holding the archive format and single-layer compression mode constant. Final conclusions should be written only after the full benchmark campaign has generated verified results.

\bibliographystyle{IEEEtran}
\bibliography{references}

\end{document}
